\section{Phương pháp thực nghiệm khoa học}

\

Để tìm ra những kết quả khoa học, chúng ta cần sử dụng \defText{phương pháp thực nghiệm khoa học}. Về mặt bằng chung, phần lớn các phương pháp đó đều phải trải qua những giai đoạn:

\begin{enumerate}[label=\protect\textbf{\color{colorEmphasis}\arabic*.}]
    \item \defText{Quan sát}: Nhận ra một hiện tượng lạ hoặc một vấn đề nảy sinh trong tự nhiên.
    \item \defText{Đặt câu hỏi}: Xác định rõ vấn đề: ``Tại sao hiện tượng này xảy ra?'', ``Nó vận hành như thế nào?''.
    \item \defText{Xây dựng giả thuyết}: Đưa ra một câu trả lời tạm thời có tính lô-gích để giải thích hiện tượng. \label{step:hypo}
    \item \defText{Dự đoán}: Suy luận ra các hệ quả nếu giả thuyết đúng: ``Nếu giả thuyết $A$ là đúng, thì khi làm được $X$, kết quả phải là $Y$''.
    \item \defText{Thực nghiệm}: Tiến hành thí nghiệm để thu thập dữ liệu, các khẳng định, các sự thật. Đây là bước đối chiếu dự đoán với thực tế.
    \item \defText{Phân tích kết quả}: Sử dụng toán học và lập luận lô-gích để xử lí dữ liệu thu được và qua đó kiểm tra kết quả thực tế có khớp với dự đoán không.
    \item \defText{Kết luận và Điều chỉnh}:
          \begin{itemize}
              \item Nếu kết quả khớp: Giả thuyết được tạm chấp nhận.
              \item Nếu kết quả sai: Giả thuyết bị bác bỏ, ta phải quay lại bước \ref{step:hypo} để xây dựng giả thuyết mới.
          \end{itemize}
\end{enumerate}

\begin{figure}[H]
    \centering
    \begin{tikzpicture}[
            node distance=0.8cm and 2cm,
            box/.style={rectangle, draw=colorEmphasisCyan, thick, fill=colorEmphasisCyan!10, minimum width=6cm, minimum height=1cm, align=center, rounded corners=3pt},
            decision/.style={diamond, aspect=2.5, draw=colorEmphasis, thick, fill=colorEmphasis!10, minimum width=3cm, minimum height=1cm, align=center},
            arrow/.style={->, >=stealth, thick, draw=black}
        ]
        % Nodes
        \node[box] (step1) {\textbf{1. Quan sát}};
        \node[box, below=of step1] (step2) {\textbf{2. Đặt câu hỏi}};
        \node[box, below=of step2] (step3) {\textbf{3. Xây dựng giả thuyết}};
        \node[box, below=of step3] (step4) {\textbf{4. Dự đoán}};
        \node[box, below=of step4] (step5) {\textbf{5. Thực nghiệm}};
        \node[box, below=of step5] (step6) {\textbf{6. Phân tích kết quả}};
        \node[decision, below=1cm of step6] (step7) {\textbf{7. Kết quả}\\\textbf{có khớp không?}};
        \node[box, draw=colorEmphasisGreen, fill=colorEmphasisGreen!10, below=1cm of step7] (accept) {\textbf{Chấp nhận giả thuyết}};

        % Arrows
        \draw[arrow] (step1) -- (step2);
        \draw[arrow] (step2) -- (step3);
        \draw[arrow] (step3) -- (step4);
        \draw[arrow] (step4) -- (step5);
        \draw[arrow] (step5) -- (step6);
        \draw[arrow] (step6) -- (step7);
        \draw[arrow] (step7) -- node[right, font=\bfseries] {Có} (accept);

        % Feedback loop
        \draw[arrow] (step7.west) -- node[above, font=\bfseries] {Không (Bác bỏ)} ++(-3.5,0) |- (step3.west);
        \draw[arrow] (accept.east) -- ++(1.5,0) |- (step1.east);
    \end{tikzpicture}
    \caption{Sơ đồ khối phương pháp thực nghiệm khoa học}
\end{figure}

Cần lưu ý rằng quá trình nghiên cứu và thực nghiệm khoa học có tính tuần hoàn. Các nhà khoa học không chỉ dừng lại ở một kết quả, mà họ tìm cách tiếp tục mở rộng và hoàn thiện nó, làm cho nền khoa học ngày càng được hệ thống hóa và chính xác hơn. Quá trình này có thể diễn ra xuyên suốt nhiều năm, thậm chí nhiều thế hệ.

Tuy rằng phương pháp tác giả vừa mô tả là vô cùng phổ biến, không phải lúc nào một phát hiện khoa học đều đi qua đầy đủ các bước này. Một số lượng lớn những kết quả đi đến từ việc thử sai hàng loạt kết hợp với phát hiện ngẫu nhiên. Xin được đưa ra một ví dụ, Xpen-xơ Xin-vơ\footnote{Spencer Silver (1941 -- 2021).}, trong quá trình nghiên cứu chất siêu dính cho ngành hàng không, đã vô tình tạo ra những ``vi cầu'' có độ dính cực thấp. Trong một thời gian dài, không ai để ý đến phát hiện này, cho đến khi Ạt-Phrai\footnote{Art Fry (Năm sinh 1931).}, một đồng nghiệp của Xpen-xơ, đã cảm thấy vô cùng phiền phức vì những mảnh giấy đánh dấu trang nhạc cứ rơi ra ngoài trong một buổi tập hát của nhà thờ. Khi đó, ông nhớ lại loại keo ``thất bại'' của Xpen-xơ và nhận ra đó chính là thứ mình cần: một loại dấu trang dính được nhưng bóc ra dễ dàng mà không làm hỏng giấy. Thế là những mẩu giấy ghi chú như ở hình \ref{fig:post_it} được ra đời. 

\begin{figure}[H]
   \centering
   \includegraphics[width=0.6\textwidth]{images/post_it.jpg}
   \caption{Bốn mẩu giấy ghi chú với màu sắc khác nhau\\
      Nguồn: https://www.pexels.com/photo/photo-of-four-colorful-sticky-notes-6991351/}
   \label{fig:post_it}
\end{figure}

Nhưng nhìn chung, đi theo con đường nào đi chăng nữa, một người học và làm khoa học đều cần phải có những phẩm chất cơ bản: sự tìm tòi, tính chính trực, sự kiên trì, và quan trọng nhất, tính khiêm nhường. Gộp chúng, một nhà khoa học cần phải có khả năng chấp nhận sai sót của bản thân và chấp nhận những kết luận mới, kể cả khi nó trái với tư duy đang sẵn có.

Phương pháp thực nghiệm khoa học chỉ có giá trị đối với những phát biểu có thể kiểm chứng được tính đúng sai. Những phát biểu khác, như những phát biểu về đạo đức, tôn giáo hay triết học, không thể được xác minh bằng các quan sát trong thế giới vật lí, cho nên, không được khoa học nghiên cứu. Đó là một trong những điểm hạn chế, khi khoa học không có khả năng trả lời tất cả mọi câu hỏi. Và do đó, lợi dụng nhu cầu cần giải đáp mọi nghi vấn của con người, một số, có thể với ý đồ xấu, sẽ cố gắng đánh lừa người khác bằng cách mượn danh khoa học để đưa ra những phát biểu nghe giống như khoa học nhưng thực tế là không có cơ sở vật chất nào. Ngắn gọn, luôn có sự xâm nhập của các thông tin \defText{giả khoa học}. Lấy ví dụ:
\begin{quotation}
    ``Tồn tại một loại năng lượng bí ẩn, không thể phát hiện ra bằng bất kì thiết bị nào, nhưng có thể chữa lành mọi bệnh tật nếu bạn biết cách sử dụng nó đúng cách.''.
\end{quotation}
Đây là một phát biểu giả khoa học điển hình. Nó có vẻ rất hấp dẫn, nhưng nó không thể được chứng minh, vì nó dựa trên một loại năng lượng không thể phát hiện được. Cho nên, nó là một mệnh đề vô giá trị. Tương tự, cũng có một giả định luật gọi là ``luật hấp dẫn'', không phải là định luật vạn vật hấp dẫn với Niu-tơn, mà thường được phát biểu là:
\begin{quotation}
    ``Bạn sẽ thu hút những điều bạn nghĩ đến. Nếu bạn nghĩ tích cực, bạn sẽ thu hút những điều tích cực; nếu bạn nghĩ tiêu cực, bạn sẽ thu hút những điều tiêu cực.''.
\end{quotation}
Giả sử tác giả là một người tuyên truyền phát biểu này, và bạn đọc muốn kiểm chứng nó. Bạn đọc có thể thử nghĩ tích cực trong một khoảng thời gian dài và xem liệu những điều tích cực có đến không. Nhưng nếu không có những điều tích cực đến, tác giả có thể đổ lỗi cho việc bạn đọc chưa nghĩ đủ tích cực, hoặc bạn đọc đã nghĩ tiêu cực vào một lúc nào đó mà bạn đọc không nhớ. Do đó, không có cách nào để bác bỏ phát biểu này, và nó trở thành một mệnh đề giả khoa học. Tệ hơn cả, những phát biểu như vậy thường được sử dụng để mê dụ những con người đang gặp khó khăn trong đời sống và đẩy họ làm những hành động trái với lí luận và đạo đức, nhất là khi mà họ chưa đầy đủ nhận thức hoặc đang không ở trạng thái tốt nhất.